\section{引言}\label{sec:intro}

\subsection{问题的提出}

图上的很多任务，人类已经写出了“最优解”：单源最短路有 Dijkstra 与 Bellman-Ford，最小生成树有 Kruskal 与 Prim。
这类算法的共同点是\textbf{白盒}的——每一步都能被审计，正确性有证明，复杂度有上界，输入被扰动之后答案会怎么变，往往也能算出来。

与之相对，图神经网络（GNN）把计算交给一层层消息传递，用数据拟合映射。
它在节点分类、链接预测上常常表现很好，但\emph{为什么}给出这个答案，只能靠事后解释器去近似~\cite{ying2019gnnexplainer}。
这类模型是\textbf{黑盒}的：假设类无法逐条读出，误差没有上界，解释是近似的近似。

题目给出的“可解释性 vs. 性能”由此成立。但本文想指出一个容易被忽略的前提：\textbf{这条权衡曲线只在算法未知时才存在}。
知道算法，就把算法写进结构里，三个轴可以同时赢；不知道算法，才需要在解释与性能之间做取舍——而且那时连“解释得对不对”都难以定义。

\subsection{为什么需要两个任务}

只用最短路一个任务，不足以支撑上面这个判断。因为在最短路里白盒是\emph{在场}的，我们只能证明“白盒在场时它赢”，不能说明“白盒缺席时会发生什么”。

所以本文做两个任务：

\begin{itemize}[leftmargin=1.6em]
\item \textbf{第一部分：算法已知}（最短路）。白盒可以精确求解，因此可以把两类方法放进\emph{同一张表}、用\emph{同一批图}直接对照。这一部分回答：当算法已知时，学习还有没有价值？
\item \textbf{第二部分：算法未知}（Cora 节点分类）。没有人能写出“给定图，输出每个节点的类”的正确程序。白盒缺席，问题从“谁更好”变成\emph{“我们还能不能验证一个解释”}。
\end{itemize}

两个任务是互补的：前者把白盒与黑盒放在同一天平上，后者则考察天平的一端不存在时怎么办。

\subsection{本文的工作}

\begin{enumerate}[leftmargin=1.8em]
\item \textbf{等价性}：$k$ 层 min-sum 消息传递与同步 Bellman-Ford 的第 $k$ 轮逐点等价（定理~\ref{thm:equiv}）。由此“深度”与“边数”被绑定，深度下界成为一个推论而不是类比。
\item \textbf{深度下界}：给出构造性反例——两张图在 $t$ 的 $L$ 跳计算图内同构（且特征逐点相同），真值却不同，因此任何 $L$ 层模型在\emph{理想参数}下也必有一个错（命题~\ref{prop:depth}）。
\item \textbf{白盒鲁棒性 = 可证明性}：删边单调与非树边不变各在 500/500 例上零违反，单边删除上界在 398 个非平凡例上零违反（另 102 例上界为 $\infty$，平凡成立）。
\item \textbf{三轴实证（第一部分）}：未训练恒等 min-sum 的误差 $\sim10^{-6}$，训练 2 个参数后退化到 $0.024$；黑盒的 Fidelity$^+$ 在 28 张图的解释集上高达 0.978，而同一集合上的相关系数只有 0.13（更难的 60 张图集上是 $-0.10$）——保真度会被常数预测器骗过。
\item \textbf{解释的可验证性（第二部分）}：在 Cora 上测攻击与解释；由于真实数据没有“正确的解释”，另造带真值的 SBM 环境，用消融定义真值，测得解释器达到 oracle 的 93.9\%；并发现攻击显著降低准确率却几乎不改变解释。
\end{enumerate}

\subsection{全文结论}

两条线合起来给出一个\emph{有条件}的回答：

\begin{quote}
当算法已知，白盒在三个轴上同时占优（在允许迭代到收敛的预算下；同深度时并非处处占优，见表~\ref{tab:same}），“学习”甚至可以是负收益；
当算法未知，白盒缺席，“可解释”本身要先被定义和验证——本文只做到这一步，\textbf{没有画出性能与解释质量之间的权衡曲线}。
\end{quote}

换句话说：本文证明了\textbf{权衡曲线在算法已知时不存在}，并说明了\textbf{在算法未知时它为什么难以被测量}——因为那一侧连“解释得对”都缺少定义，只能用合成环境去近似。
至于曲线在未知一侧的具体形状，是一个未解决的问题；第~\ref{sec:discussion}~节会说明为什么它比看上去更难测。

\subsection{全文结构}

第~\ref{sec:white}~节给出白盒算法与轮数语义；第~\ref{sec:robust}~节把鲁棒性写成可证命题；第~\ref{sec:black}~节建立消息传递与 Bellman-Ford 的等价并给出深度下界；
第~\ref{sec:resultsA}~节报告第一部分的实验；第~\ref{sec:node}~节转向节点分类，说明没有白盒时怎么做；第~\ref{sec:noderes}~节报告第二部分的攻击与解释结果；
第~\ref{sec:discussion}~节综合两条线；第~\ref{sec:conclusion}~节总结。
